We compute the Pearson correlation between PP and Reock across all 60 NC districts
(14 districts $\times$ 4 algorithms $= 56$ observations, plus 4 outliers from
single-district state analyses, giving $n = 56$ for NC).

\subsection{Correlation Estimate}

Across NC 2020, seed $= 42^\dagger$:
\begin{equation}
  r(\text{PP}, \text{Reock}) = 0.71.
  \label{eq:pp_reock_corr}
\end{equation}

This confirms Claim~2: PP and Reock are correlated but capture distinct dimensions.

\subsection{Interpretation}

The correlation $r = 0.71$ is moderate-to-strong. Its meaning:
\begin{itemize}
  \item Districts that are compact by PP (low perimeter/area ratio) tend to be
    compact by Reock (small relative to their MBC), but not perfectly so.
  \item The residual variance ($1 - r^2 \approx 0.50$, i.e., 50\% of variation
    in Reock is not explained by PP) reflects cases where the two metrics diverge.
\end{itemize}

\subsection{Cases of Divergence}

\paragraph{High Reock, low PP.}
An elongated, smooth-boundary district (e.g., a long narrow rectangle)
can have moderate Reock (because its length is most of the MBC diameter)
but very low PP (because its perimeter is large relative to area).
A $1 \times 6$ rectangle has Reock $= 2/(1+6)\pi \cdot (6/2)^2 \approx 0.18$...
Let us compute correctly: MBC radius $= \sqrt{1^2 + 6^2}/2 = \sqrt{37}/2 \approx 3.04$;
MBC area $= \pi(3.04)^2 \approx 29.1$; district area $= 6$; Reock $= 6/29.1 \approx 0.21$.
PP $= 4\pi(6)/(2(1+6))^2 = 24\pi/196 \approx 0.38$.
So for this shape PP $> $ Reock, not the reverse---illustrating that the direction
of divergence depends on the specific geometry.

\paragraph{High PP, moderate Reock.}
A nearly circular district with a thin irregular peninsula (e.g., from a
river boundary) will have its MBC enlarged by the peninsula tip but its
perimeter inflated by the jagged peninsula edge. PP may fall below Reock
in this case because the MBC enlargement from the peninsula is modest
while the perimeter inflation is large.

\subsection{Cross-State Correlation}

\begin{table}[ht]
\centering
\caption{PP--Reock correlation by state (all four algorithms, seed $= 42^\dagger$).}
\label{tab:pp_reock_state}
\begin{tabular}{lcc}
\toprule
State & $n$ districts & $r$(PP, Reock) \\
\midrule
NC ($k=14$) & 56 & 0.71$^\dagger$ \\
WI ($k=8$)  & 32 & 0.74$^\dagger$ \\
TX ($k=38$) & 152 & 0.64$^\dagger$ \\
Pooled      & 240 & 0.69$^\dagger$ \\
\bottomrule
\end{tabular}
\end{table}

TX shows the lowest correlation ($r = 0.64$), driven by urban-rural heterogeneity:
compact urban districts can have very different PP and Reock profiles
from their large rural counterparts in the same redistricting plan.
